% =====================================================================
% Rede textual da tese lida pelos capítulos — seção de síntese
% Modelo: esboco_secao_infranodus.tex (capítulo 1).
%
% Esta seção NÃO descreve o grafo da capa do site (tese_network.py), que
% roda um único Louvain sobre a tese concatenada. Ela lê comparativamente
% as quatro redes de capítulo (figuras \ref{fig:infranodus-cap1} a
% \ref{fig:infranodus-cap4}), rastreando o que recorre e o que se desloca
% de um capítulo a outro.
%
% Estilo seguido (regras da tese): primeira pessoa do singular, sem
% travessões em texto principal, "capítulo" minúsculo, aspas com
% \enquote{}, sem fórmulas do tipo "não X, mas Y".
% =====================================================================

\subsection*{A tese lida pela sua própria rede: uma trajetória entre capítulos}

Reunidas as quatro redes de capítulo (figuras~\ref{fig:infranodus-cap1}
a~\ref{fig:infranodus-cap4}), passo a lê-las em série, atento ao que
recorre e ao que se desloca. Três achados atravessam os quatro grafos e
compõem, tomados juntos, uma cartografia da própria tese.

O primeiro achado é uma espinha lexical estável. O par
\textit{inteligência}~$\leftrightarrow$~\textit{artificial} é a
co-ocorrência de associação mais alta em três dos quatro capítulos
(NPMI~0,83 no capítulo~\ref{capitulo2}, 0,86 no
capítulo~\ref{capitulo3}, e forte já no capítulo~\ref{capitulo1}): o
objeto desta tese aparece nomeado sempre da mesma forma, colado, em todo
o percurso. Ao lado dele, os termos \textit{rede} e \textit{pesquisa}
figuram entre os de maior \textit{degree} e maior intermediação em todos
os capítulos, o que os torna as palavras-coringa que conduzem o trânsito
entre os vocabulários. A espinha do texto é, assim, feita de poucos
termos transversais, enquanto os vocabulários próprios de cada capítulo
(o \textit{militar} da genealogia no capítulo~\ref{capitulo2}, o
\textit{hollerith} e a \textit{tabulação} da infraestrutura no
capítulo~\ref{capitulo3}, o \textit{spira} e o \textit{covideiro} da cena
clínica no capítulo~\ref{capitulo4}) permanecem locais, ancorando cada
etapa do argumento no seu campo.

O segundo achado é a autonomia lexical do método. Em cada um dos quatro
capítulos a detecção de comunidades isola um agrupamento próprio da
teoria ator-rede, reunido em \textit{rede}, \textit{ator},
\textit{associação}, \textit{latour} e, no capítulo~\ref{capitulo3},
nos verbos em primeira pessoa \textit{sigo} e \textit{descrevo}. O
vocabulário com que descrevo tem existência própria e viaja pela tese
inteira. A recorrência tem, porém, uma contraface que só a leitura em
série revela, e é o terceiro achado: em todos os capítulos a ligação
ponderada mais fraca cai entre esse vocabulário do método e o vocabulário
do objeto empírico. É a fissura entre \textit{método} e
\textit{artificial} no capítulo~\ref{capitulo1}, entre a leitura das
figurações e o panorama bibliométrico no capítulo~\ref{capitulo2}, entre
o método ator-rede e o centro de pesquisa no capítulo~\ref{capitulo3}, e
entre a teoria da inscrição e o par \textit{dado}~$\leftrightarrow$~%
\textit{modelo} no capítulo~\ref{capitulo4}. A rede devolve, capítulo
após capítulo, a mesma tarefa: costurar com mais força o vocabulário com
que prometo descrever ao vocabulário daquilo que descrevo.

Sobre essa recorrência desenha-se um movimento, e é o que me autoriza a
falar de trajetória. A figura da inscrição migra ao longo da tese. Ela
está latente no capítulo~\ref{capitulo1}, apenas esboçada num pequeno
agrupamento; ganha o estatuto de objeto de leitura no
capítulo~\ref{capitulo2}, quando trato das figurações e da lexicometria
das obras de Latour; e converte-se, no capítulo~\ref{capitulo4}, em nó
central do grafo, colada ao SPIRA nos pares
\textit{imutável}~$\leftrightarrow$~\textit{móvel} e
\textit{caixa}~$\leftrightarrow$~\textit{preta}. Lida pela sua própria
rede, a tese aparece como a lexicalização progressiva do aparato
latouriano da inscrição sobre um objeto técnico determinado: o que
começa como vocabulário teórico distante termina como descrição colada à
matéria acústica e estatística de um sistema de inteligência artificial.

Registro, por fim, uma ausência que a série torna eloquente. A
contrafiguração têxtil-feminista que oponho, no capítulo~\ref{capitulo2},
à figuração militar-industrial não deixa rastro entre os termos centrais
de nenhum dos quatro grafos, ao passo que \textit{militar},
\textit{inscrição} e \textit{rede} se inscrevem com nitidez. A
desproporção mede, na letra da própria tese, o desafio que ela enuncia,
já que a figuração dominante está sedimentada no vocabulário e a
alternativa que proponho ainda precisa ser tecida no texto para deixar
traço. A rede textual inscreve esse desequilíbrio sem julgá-lo, e é
essa inscrição que reconheço como o retrato mais honesto do que escrevi,
do que prometi e do que ainda me falta escrever \parencite[p.~143]{Law2004After}.
